% !TeX program = xelatex
\documentclass{jingyan}

\begin{document}

\ArticleTitle{文章标题}{作者姓名}
\ArticleAbstract{此处填写摘要，简述研究问题、主要论证和结论。}
\ArticleKeywords{关键词一；关键词二}

\section{问题的提出}
在此填写正文。每个自然段之间空一行。

\section{分析与论证}
\subsection{概念与材料}
叙述式引用示例为\TextCite{作者甲}{2025}[第34页]，括号式引用示例为\ParenCite{作者甲}{2025}。

\subsection{公式示例}
\begin{equation}\label{eq:mean}
  \bar{x}=\frac{1}{n}\sum_{i=1}^{n}x_i.
\end{equation}
式\eqref{eq:mean}展示了公式编号和交叉引用的写法。

\subsection{表格示例}
表~\ref{tab:procurement}列出两类采购的示例金额和占比。

\begin{table}[htbp]
  \centering
  \small
  \setlength{\belowcaptionskip}{4pt}
  \caption{采购构成示例}
  \label{tab:procurement}
  \begin{tabular}{lrr}
    \toprule
    采购类别 & 采购额（万元） & 占比（\%） \\
    \midrule
    本地采购 & 60 & 60 \\
    外地采购 & 40 & 40 \\
    \bottomrule
  \end{tabular}
\end{table}

\section{结论}
在此填写结论。

\begin{References}
  \Reference{作者甲，2025：《文章题名》，《刊物名称》第X期，第XX—XX页。}
\end{References}

\end{document}
